\worldchapter{The Character}{character}
\label{ch:character}

\begin{multicols}{2}

To enter the world of \trans{book_title}, you need a character - a fictional person you breathe life into and whose destiny you guide. The character sheet is where all the stats and traits that make your hero unique are recorded.

This chapter provides an overview of a character's fundamental attributes. The process of character creation is described in detail in \cref{ch:create}.

\section{Persona}

The Persona attributes refer to the mental characteristics of the character. Each value corresponds to a personality trait. Persona traits have their own value and are also the base for skills.

\subsection{Education}

Education measures a character's acquired general knowledge and their ability to recall learned information. A high value suggests a person who has read a lot, attended a good school, or possesses an insatiable curiosity. Consequently, such a character excels in all theoretical skills like \textit{Nature} or \textit{History}.

\subsection{Logic}

While \textit{Education} represents a character's stored knowledge, Logic is their ability to apply that knowledge and draw new conclusions. Logic is always crucial when it comes to creating a coherent picture from existing clues or analyzing complex problems. A character with high Logic is therefore superior in skills like \textit{Investigation}, to connect the crucial details at a crime scene, or in \textit{Mechanics}, to see through the design of a trap and disarm it.

\subsection{Conscientiousness}

Conscientiousness describes a character's diligence, discipline, and reliability. A high value signifies a methodical approach and prevents careless mistakes, which is crucial for skills like \textit{First Aid} or \textit{Stealth}.

\subsection{Willpower}

Willpower is the mental fortitude and determination with which a character pursues their goals and resists adversity. It represents the inner toughness to not be swayed from one's path by external pressure or internal doubt. A high value is therefore the foundation for skills like \textit{Courage}, to remain steadfast even in hopeless situations, or \textit{Intimidation}, to project one's own will onto others.

\subsection{Apprehension}

Apprehension describes how quickly and precisely a character grasps and processes their surroundings with all senses. It is the measure of innate alertness and awareness of one's environment, from the smallest sound to the widest landscape. A high value is therefore the foundation for skills like \textit{Perception}, to spot hidden details or dangers, and \textit{Orientation}, to avoid getting lost in unfamiliar terrain.

\subsection{Charm}

Charm is the ability to create a positive connection with others and win them over through choice of words, demeanor, and personality. In contrast to purely external \textit{Attractiveness}, Charm is an attribute based on intuition. A high value in this area is the foundation for skills like \textit{Politics}, to win allies and conduct negotiations, as well as \textit{Empathy}, to understand and react to the moods of others.

\section{Physis}

All physical attributes describe the character's physical abilities. Each attribute has a value that indicates the number of dice rolled for that attribute.

\subsection{Deftness}

Deftness covers both a character's general physical control as well as their fine hand-eye coordination. It is the measure of reflexes, balance, and the ability to execute precise and controlled movements. A high value in this attribute allows a character to excel in skills like \textit{Acrobatics}, to evade obstacles and maintain balance, as well as in \textit{Shooting}, to reliably hit a target even at a great distance.

\subsection{Strength}

Strength is the measure of a character's raw muscle power and physical might. It represents the ability to exert overwhelming physical force, whether to move heavy objects or to inflict devastating damage in combat. A high Strength value is therefore crucial for \textit{Hand to hand combat}, to smash through armor with powerful blows, as well as for \textit{Throwing}, to hurl objects at a target with great force.

\subsection{Attractiveness}

Attractiveness measures the immediate impact of a character's physical appearance and presence on others. This value doesn't necessarily describe conventional beauty, but rather how memorable or captivating a person's appearance is - be it through graceful features, an intimidating stature, or striking characteristics. High Attractiveness ensures that a character stands out from the crowd and leaves a strong first impression before they have even spoken a word.

\subsection{Endurance}

Endurance describes a character's purely physical resilience and toughness. It determines how long someone can engage in strenuous activities like a forced march, a long run, or a fight before exhaustion sets in.

\subsection{Resistance}

Resistance is the body's innate toughness and constitution, allowing a character to withstand damage and hostile influences. It represents the ability to mitigate the effects of poisons, or endure the hardships of extreme heat and cold. A character with high Resistance fends off diseases and other harmful effects more effectively.

\subsection{Quickness}

Quickness measures both a character's pure movement velocity as well as their reaction time. It determines how rapidly a person can travel to cover distances, but also how quickly they can react to sudden events or dangers.

\section{The minimum roll}

The minimum roll is a central feature of the character. It specifies the result a die must have to represent a success. The minimum roll is defined by the lineage and is 5+ for most characters. The ``Masterly Presence'' template lowers the minimum roll by 1, otherwise it can only be changed by special events or rare items, and often only for a short time.

\section{Evasion}

\textit{Evasion} is used in combat and allows you to avoid a melee attack. It is equal to the Evasion value of the lineage plus the average of \textit{Quickness} and \textit{Deftness} (rounded up). Armour and weapons reduce this stat. Character templates can change this value.

\section{Protection}

If a character has protection due to their lineage, this is called ``innate protection''. Unlike the other protection types, this protection does not deplete until the end of combat; it refreshes at the start of the player's combat round (see \cref{ch:combat}).

\section{Additional dice}

Each character can have a number of \textit{bonus dice}, \textit{destiny dice} or \textit{rerolls}. All three have different uses (see \cref{ch:rolls}), but always represent an advantage to the character that can be used during the game.

The character can regain used dice during the rest (see \cref{ch:wounds}).

\section{Skills}

More complex actions or knowledge are described by \textit{skills}. All characters have the same skills with different values, so the GM can be sure that a player can definitely roll a skill.

Each skill has a base attribute and a skill value. For example, the base attribute for the skill \textit{Intimidation} is \textit{Apprehension}.

The base attribute is added to the bonuses of the selected character templates.

\subsection{Intimidation}

The Intimidation skill is a measure of how well a character can intimidate others. This skill can be used to extract information from an opponent or to make them retreat from a fight.

Attribute: \textit{Apprehension}

\subsection{Empathy}

Empathy is the ability to interpret a person's feelings and moods, and perhaps to recognise intentions. Thoughts cannot be read.

Attribute: \textit{Charm}

\subsection{Stealth}

Stealth is the art of concealment. This skill is used both for stealthy movement (sneaking) and to check how secretive the character is. It can be used, for example, when the character is being interrogated or is tempted to divulge a secret.

Attribute: \textit{Conscientiousness}

\wbif{nexus}{}{

\subsection{Magic Knowledge}

\wbif{phasesix}{This skill is only available for campaigns containing magic.}{}

Magic Knowledge is the theoretical and academic knowledge about the nature of magic. It allows a character to analyze the structure of a foreign spell, decipher arcane writings, or know the history and rituals of ancient magical cults. In direct contrast to the practical skill of \textit{Spell Casting}, Magic Knowledge is the skill of understanding and identifying, not of applying.

Attribute: \textit{Charm}
}

\subsection{Orientation}

This skill is used for orientation, both in the countryside and in confusing situations. It can be used in the confusing crowds of the city, but also when the character is whirled by a water vortex.

Attribute: \textit{Apprehension}

\subsection{Politics}

Whenever it comes to assessing political action, this skill is used.
\wbif{tirakan}{This can be the case in real politics, but can also represent moving safely in aristocratic circles.}{This can be the case in real politics, but can also represent moving safely in large corporations.}

Attribute: \textit{Charm}

\subsection{Religion}

This skill includes knowledge of religious teachings, as well as confidence in performing religious ceremonies.

Attribute: \textit{Conscientiousness}

\subsection{Courage}

This skill comes into play whenever it is a question of how brave a character is. For example, it can be used to determine whether a character is brave enough to face a powerful opponent.

Attribute: \textit{Willpower}

\subsection{Deception}

If the character wants to deceive an opponent, or, for example, cheat at the game, this skill can be rolled on.

Attribute: \textit{Charm}

\subsection{Persuasion}

If the character wants to convince his counterpart argumentatively, this skill is used.

Attribute: \textit{Willpower}

\subsection{Investigation}

This skill is used when the character wants to examine an object, a certain scene or an object for certain properties.

Attribute: \textit{Apprehension}

\subsection{Perception}

Perception represents the character's ability to perceive things in his environment. This can be the search of an house, the search for the shadowy thief at the edge of the forest, or even a movement in the face of the opponent.

Attribute: \textit{Apprehension}

\subsection{Acrobatics}

Acrobatics is the art of moving quickly and skillfully. Unlike athleticism, this skill is used when the character climbs over a ledge or makes a short sprint.

Attribute: \textit{Deftness}

\subsection{Performance}

Performance is the artistic presentation. This can be acting, but also the musical performance of a piece. An impressive tall tale can also be told with the help of performance.

Attribute: \textit{Charm}

\subsection{First Aid}

First aid must be carried out with sufficient dressing materials to be successful.

If the throw is successful, the person receiving first aid recovers wounds equal to half the successes (rounded up) of the throw.

First aid stops any bleeding.

Attribute: \textit{Conscientiousness}

\subsection{Driving}

The Driving skill describes the driving of all kinds of vehicles. The skill applies to all mobile objects such as ships, vehicles or carriages.

Attribute: \textit{Deftness}

\subsection{History}

History describes the character's knowledge of history and past events. Antiquities can also be assessed with this skill.

Attribute: \textit{Education}

\subsection{Communication}

The ability to socialize is described by the skill Communication. It describes how skillfully the character behaves in conversations.

Attribute: \textit{Education}

\subsection{Mechanics}

Mechanics includes all manual activities as well as the knowledge of mechanical processes. Working on a piece of wood or understanding a mechanical clock can be mapped with this skill.

Attribute: \textit{Logic}

\subsection{Hand To Hand Combat}

The value of this skill is the basis for attacking with melee weapons. This skill is not usually rolled on directly.

Attribute: \textit{Strength}

\subsection{Nature}

Nature describes the character's knowledge of all facets of nature. This skill can be used when the character is searching for plants, gathering wood in the forest, or judging the nature of an animal.

Attribute: \textit{Education}

\subsection{Shooting}

The value of this skill is the basis for attacking with ranged weapons. This skill is not usually rolled on directly.

Attribute: \textit{Deftness}

\subsection{Throwing}

This skill is used whenever the character throws objects. These can be simple objects like stones, but also incendiary charges or nets.

For exact rules on throwing items, see \cref{sec:throwing}.

Attribute: \textit{Strength}

\wbif{nexus}{}{
\subsection{Spell Casting}

\wbif{phasesix}{This skill is only available for campaigns containing magic.}{}

The spell casting skill is the measure of a character's mastery over raw magical energies and their ability to shape and direct them with purpose. It is used to cast powerful spells, engage in magical duels, or create artifacts. In crucial contrast to the purely theoretical \textit{Magic Knowledge}, spell casting thus describes the active and practical application of magic.

Attribute: \textit{Willpower}
}

\section{Knowledge}

Knowledge works in a similar way to skills, but the list is not predefined. Characters can have different knowledge skills based on their background, which they can use freely. Knowledge is always associated with a skill. The effective die roll value is the sum of the knowledge value and the skill value.

Knowledge is gained through character templates. The character templates indicate whether they bring this knowledge with them.

\section{Shadows}

A character can have special traits that affect them outside of their physical or mental attributes. Each \textit{shadow} has its own description or rule. For example, a character may have a rival or be obedient to authority. Shadows do not have values, but can have their own rules.

Shadows are indicated on character templates. If a character template contains a written rule, it is a shadow.

\section{Languages}

The number of languages a character can learn is based on the sum of their \textit{Education} and \textit{Logic} attributes. These can be any languages from the character's world. If the sum of these attributes is 0 or less, the character has only a limited understanding of their native language.

The limit on the number of languages that can be learnt serves as a guideline for new characters. However, languages learned in the course of the game can exceed this limit.

\wbif{nexus}{
Character templates or body modifications can increase the number of languages that can be learnt.
}{}
\wbif{tirakan}{
Character templates or magical items can increase the number of languages that can be learnt.
}{}
\wbif{phasesix}{
Character templates, body modifications or magical items can increase the number of languages that can be learnt.
}{}

\section{Contacts}

Contacts are connections that a character has with other people or beings that they can rely on. These are typically people outside the party, such as a noble, a government contact, or a doctor.

When creating a character, they can have a certain number of contacts, based on the sum of the \textit{Charm} and \textit{Attractiveness} attributes.

This number can be exceeded if new contacts are made during the game.
\end{multicols}
